\documentclass[exercice,a4]{thedocument}%
\usepackage[french]{babel}%
\usepackage{tabularx}%
\usepackage{array}%
\startdatakeys
    \source{21GENMATNC1}
    \BO{2026}
    \cycle{4}
    \competencesDuSocle{CA,RA}
    \objectifsApprentissage{C4C1ac,C4C2aa,C4C2ab}
    \motsCles{moyenne,étendue,médiane,pourcentage}
\enddatakeys
\begin{document}%

\medskip

On étudie les précipitations (hauteurs de pluies) sur la ville de Nouméa entre avril et décembre 2020.

On obtient le tableau suivant:

\begin{center}
\begin{tabularx}{\linewidth}{|m{3.5cm}|*{9}{>{\centering \arraybackslash}X|}}\hline
Mois 				&Avril 	&Mai &Juin 	&Juillet &Août 	&Sept 	&Oct &Nov 	&Déc.\\ \hline
Précipitations en mm&147 	&199 &40 	&67 	&47 	&54 	&104 &45 	&63\\ \hline
\multicolumn{10}{r}{\footnotesize Source: https://www.historigue-meteo.net/oceanie/nouvelle-caledonie/noumea/2020}\\
\end{tabularx}
\end{center}

\medskip

\begin{enumerate}
\item Calculer la moyenne des précipitations. Arrondir le résultat au mm près.
\item Quelle est l'étendue des précipitations ?
\item Déterminer la médiane des précipitations.
\item Calculer le pourcentage de mois pour lesquels les précipitations sont supérieures à $100$~mm. Arrondir le résultat à l'unité près.
\end{enumerate}

\bigskip

\showthesource
\end{document}
